\documentclass[12pt,a4paper]{article}
\usepackage[T2A]{fontenc}
\usepackage[utf8]{inputenc}
\usepackage[russian]{babel}
\usepackage{amsmath}
\usepackage{amsfonts}
\usepackage{amssymb}
\usepackage{amsthm}
\usepackage{geometry}
\usepackage{fancyhdr}
\usepackage{titlesec}
\usepackage{graphicx}
\usepackage{listings}
\usepackage{float}
\usepackage{booktabs}
\usepackage{array}
\usepackage{xcolor}
\usepackage{parskip}
\usepackage{setspace}
\usepackage{microtype}
\usepackage{hyperref}

\usepackage{chngcntr} % для \counterwithin
\setcounter{tocdepth}{3} % включить в оглавление sections (1), subsections (2) и subsubsections (3)
\counterwithin{figure}{section} % нумерация рисунков будет вида <section>.<figure>

% Настройка страницы
\geometry{a4paper, left=2.5cm, right=2.5cm, top=2.5cm, bottom=2.5cm}
\setstretch{1.25} % Межстрочный интервал
\setlength{\parskip}{0.8em} % Отступ между параграфами

% Настройка стиля заголовков
\titleformat{\section}
{\normalfont\Large\bfseries\centering}{\thesection}{1em}{}
\titleformat{\subsection}
{\normalfont\large\bfseries}{\thesubsection}{1em}{}
\titleformat{\subsubsection}
{\normalfont\normalsize\bfseries\itshape}{\thesubsubsection}{1em}{}

\title{МОСКОВСКИЙ АВИАЦИОННЫЙ ИНСТИТУТ

(НАЦИОНАЛЬНЫЙ ИССЛЕДОВАТЕЛЬСКИЙ УНИВЕРСИТЕТ)

Институт №8 «Компьютерные науки и прикладная математика»

\hfill\break
\hfill\break
\hfill\break

\textbf{Лабораторные работа}

\textbf{по курсу «\emph{Информационный поиск}»}

\hfill\break

Выполнил: \emph{Потапов Егор Дмитриевич}

Группа: \emph{М8О-409Б-22}

Преподаватель: \emph{Кухтичев Антон Алексеевич}

\hfill\break

Москва, \emph{2025}}
\author{}
\date{}

\begin{document}
\maketitle

\textbf{ОБЩЕЕ ОПИСАНИЕ ПРОЕКТА}

В работе реализована цепочка компонентов поисковой системы: подготовка
корпуса HTML-документов, поисковый робот для переобкачки и устойчивого
продолжения, токенизация и статистика, анализ закона Ципфа,
морфологическая нормализация (стемминг), построение булева индекса в
собственном бинарном формате и булев поиск. Код в отчёте не приводится
(по требованию), приводятся описание решений, инструкции по запуску и
демонстрация работы.

\textbf{Используемые источники корпуса}

Источник 1: koolinar.ru (рецепты)\\
Источник 2: povarenok.ru (рецепты)

Примеры документов:\\
• https://www.koolinar.ru/recipe/view/150000\\
• https://www.povarenok.ru/recipes/show/20000/

\textbf{Важное ограничение для лабораторных работ 3--7}

Основные компоненты (индексация/поиск) реализованы на C++ без
использования структур хранения ключей (unordered\_map/map/set и
аналогов). Разрешены std::string и std::vector. Для токенизации
допускается использование STL по условиям задания.

\textbf{Лабораторная работа №1. Добыча корпуса документов}

\textbf{1. Описание корпуса и структуры документов}

Тип документов: HTML-страницы рецептов. Язык: русский. «Сырой» HTML
помимо основного контента содержит сервисные блоки (меню, рекомендации,
формы входа, рекламу, скрипты).

\begin{itemize}
\item
  Основной контент рецепта: заголовок, ингредиенты, шаги приготовления,
  теги/категории, иногда комментарии и автор.
\item
  Служебные блоки: меню, «похожие рецепты», рекомендации, футер,
  баннеры, аналитика и JavaScript.
\end{itemize}

\textbf{2. Выделение текста из HTML}

Текст выделяется из DOM-дерева следующим образом: из HTML удаляются теги
script/style/noscript, а также HTML-комментарии; далее берётся видимый
текст страницы (get\_text) и нормализуются пробелы.

\textbf{3. Проверка пригодности корпуса: наличие существующего поиска}

Требование задания: у источников должен существовать поисковый
интерфейс, чтобы можно было проверить ответы вручную.

\begin{itemize}
\item
  koolinar.ru: поиск по сайту по URL вида
  https://www.koolinar.ru/list?search\_term\_string=курица
\item
  povarenok.ru: поиск по сайту по URL вида
  https://www.povarenok.ru/recipes/poisk/kurica/
\end{itemize}

\textbf{4. Примеры запросов к существующему поиску и недостатки выдачи}

\begin{itemize}
\item
  Запрос «курица» (оба сайта): много результатов, трудно выделить точные
  совпадения по составу блюда.
\item
  Запросы со словами в разных формах («помидоры»/«помидором»): качество
  зависит от морфологии и настроек поиска сайта.
\item
  Нет возможности задать строгий булев запрос вида (курица \&\&
  помидоры) \&\& !майонез.
\end{itemize}

\textbf{5. Статистика по корпусу}

Ниже приведён пример статистики для 40000 документов.

\begin{longtable}[]{@{}
  >{\raggedright\arraybackslash}p{(\columnwidth - 2\tabcolsep) * \real{0.5000}}
  >{\raggedright\arraybackslash}p{(\columnwidth - 2\tabcolsep) * \real{0.5000}}@{}}
\toprule
\begin{minipage}[b]{\linewidth}\raggedright
Показатель
\end{minipage} & \begin{minipage}[b]{\linewidth}\raggedright
Оценка 40 000 док.
\end{minipage} \\
\begin{minipage}[b]{\linewidth}\raggedright
Кол-во документов
\end{minipage} & \begin{minipage}[b]{\linewidth}\raggedright
40000
\end{minipage} \\
\begin{minipage}[b]{\linewidth}\raggedright
Raw объём, байт
\end{minipage} & \begin{minipage}[b]{\linewidth}\raggedright
2 061 822 160
\end{minipage} \\
\begin{minipage}[b]{\linewidth}\raggedright
Текстовый объём, байт
\end{minipage} & \begin{minipage}[b]{\linewidth}\raggedright
314 781 120
\end{minipage} \\
\begin{minipage}[b]{\linewidth}\raggedright
Кол-во слов
\end{minipage} & \begin{minipage}[b]{\linewidth}\raggedright
25 234 640
\end{minipage} \\
\begin{minipage}[b]{\linewidth}\raggedright
Средний raw, байт/док
\end{minipage} & \begin{minipage}[b]{\linewidth}\raggedright
≈ 51 545.6
\end{minipage} \\
\begin{minipage}[b]{\linewidth}\raggedright
Средний текст, байт/док
\end{minipage} & \begin{minipage}[b]{\linewidth}\raggedright
≈ 7 869.5
\end{minipage} \\
\begin{minipage}[b]{\linewidth}\raggedright
Среднее кол-во слов/док
\end{minipage} & \begin{minipage}[b]{\linewidth}\raggedright
≈ 630.9
\end{minipage} \\
\midrule
\endhead
\bottomrule
\end{longtable}

\textbf{Инструкция по запуску (ЛР1) и демонстрация}

Сбор HTML и статистики выполняется скриптом parse.py (пример).
Результатом является директория recipes\_corpus/ с поддиректориями raw/,
text/, meta/.

cd /home/egor/search/1

python3 parse.py

ls -la recipes\_corpus/povarenok/meta/meta.jsonl \textbar{} head

cat recipes\_corpus/povarenok/povarenok\_stats.json

\textbf{Лабораторная работа №2. Поисковый робот}

Цель работы --- реализовать поискового робота, который обходит страницы
по заданным правилам и сохраняет результаты обкачки в MongoDB. В отличие
от ЛР1, промежуточная запись в файлы не используется: документ сразу
сохраняется в БД в требуемом формате.

Входной параметр программы --- путь до YAML-конфига. Конфиг содержит
секции: db (подключение к MongoDB и имена коллекций), logic (паузы между
запросами, таймауты, интервалы переобкачки/повторов), headers
(User-Agent), collect (правила генерации URL по шаблонам и целевое число
документов по каждому источнику).

Используются три коллекции MongoDB:

• urls --- очередь URL для обхода (url, source, next\_crawl\_ts,
last\_crawl\_ts, status\_code, hash). next\_crawl\_ts обеспечивает
планирование следующей попытки/переобкачки.

• docs --- хранилище документов (url --- нормализованный, raw\_html ---
«сырой» HTML, source --- название источника, crawl\_ts --- дата обкачки
в формате Unix timestamp). Дополнительно сохраняются title и text,
извлечённые из HTML (понадобятся в последующих лабораторных работах).

• collect\_state --- состояние генератора URL по источникам (site\_key,
next\_id), что позволяет продолжать работу после остановки.

Нормализация URL: приводится схема и хост к нижнему регистру, удаляется
фрагмент (\#...), путь очищается от лишних «//» и нормализуются
завершающие слэши.

Механизм «переобкачки только если документ изменился» реализован
сравниванием хеша контента: после загрузки вычисляется md5(raw\_html) и
сравнивается с сохранённым значением в urls.hash. Если хеш изменился ---
документ в docs обновляется, иначе обновляется только информация о
времени обхода и планирование следующей проверки.

Остановка и возобновление: обработан сигнал SIGINT/SIGTERM --- программа
завершает цикл после текущей итерации; при повторном запуске состояние
очереди urls и collect\_state сохраняется в MongoDB, поэтому обход
продолжается без потери прогресса.

\textbf{Инструкция по запуску (ЛР2)}

docker run -d -\/-name mongo -p 27017:27017 -v mongo\_data:/data/db
mongo:7

docker ps

pip install pyyaml requests pymongo beautifulsoup4

python crawler.py config.yaml

mongosh

Внутри

"""use crawler\_db

db.docs.countDocuments()

db.docs.find(\{\}, \{url:1, source:1, crawl\_ts:1,
title:1\}).limit(3).pretty()"""

\textbf{Лабораторная работа №3. Токенизация}

Цель работы --- реализовать токенизацию корпуса и собрать статистику:
число токенов, среднюю длину токена, время обработки и скорость
токенизации.

В данной версии пайплайна (после доработки ЛР2) входные документы
берутся напрямую из MongoDB из коллекции docs, поле text. Это исключает
необходимость перегонять документы в *.txt и ускоряет итерации (нет
файлового I/O и проблем с путями).

Конфигурация хранится в YAML. Секция db задаёт uri, database и
docs\_collection; секция corpus задаёт max\_docs; секция tokenizer
задаёт правила токенизации (lowercase, normalize\_yo, keep\_numbers,
min\_len).

Загрузка корпуса реализована модулем mongo\_corpus
(mongo\_corpus.h/.cpp): LoadMongoConfigFromYaml читает параметры
подключения, LoadDocsFromMongo загружает до max\_docs документов,
возвращая массив структур \{doc\_id, url, title, text\}.

Алгоритм токенизации:

1) при необходимости приводим к нижнему регистру (lowercase);

2) нормализуем «ё→е» (normalize\_yo);

3) разбиваем текст на токены по неалфавитно-цифровым разделителям;

4) фильтруем токены по минимальной длине (min\_len) и (опционально) по
признаку «число» (keep\_numbers).

В конце вычисляются метрики: total\_tokens, avg\_token\_len\_bytes,
time\_sec и speed\_kb\_per\_sec (скорость в КБ/сек по суммарному размеру
текстов).

\textbf{Результаты и выводы (ЛР3)}

Эксперимент : 40000 документов, общий объём текста 1309376.230 КБ.

Токенов: 108 217 000, средняя длина токена: 10.8474 байт.

Время: 22.8400 с; скорость: 57343.6 КБ/с (\textasciitilde56 МБ/с).

\textbf{Инструкция по запуску (ЛР3)}

cd /home/egor/search/345

mkdir -p build

cd build

cmake .. \&\& make -j

./lab3\_token\_stats ../config.yaml

\textbf{Лабораторная работа №4. Закон Ципфа}

Цель работы --- построить распределение частот терминов и сравнить его с
законом Ципфа в логарифмической шкале.

Документы загружаются из MongoDB аналогично ЛР3: используется коллекция
docs, поле text. Токенизация выполняется тем же модулем Tokenizer с
настройками из YAML.

Далее все токены собираются в один массив, сортируются собственной
сортировкой (merge sort), после чего агрегируются частоты (частота терма
= длина группы одинаковых токенов в отсортированном массиве).

Частоты сортируются по убыванию (также merge sort). Для ранга r и
частоты f(r) строится теоретическая кривая Ципфа z(r)=C/r, где C=f(1).
Результат сохраняется в zipf.csv с колонками rank,freq,zipf.

\textbf{Результаты и выводы (ЛР4)}

\includegraphics[width=6.72847in,height=5.04167in]{media/image1.png}

Время построения частот и формирования zipf.csv: 164.43 с.

Оценка для 40 000 документов: токенов \textasciitilde108 217 020 (≈108.2
млн).

Если считать по измеренной производительности \textasciitilde{}
658 134,282065 токен/с, время \textasciitilde{} 2,7333 мин.

Вывод по закону Ципфа: на малой выборке «хвост» распределения обрезан
(мало редких термов), поэтому совпадение с 1/r хуже.

Оптимальность и ускорение: текущая скорость уже высока для однопоточной
обработки. Ускорить можно за счёт:

\begin{itemize}
\item
  \begin{quote}
  потоковой обработки документов (несколько потоков по документам),
  \end{quote}
\item
  \begin{quote}
  уменьшения аллокаций (переиспользование буферов tokens, reserve),
  \end{quote}
\item
  \begin{quote}
  более быстрого прохода по UTF-8 (сканирование байтами + таблицы
  классов),
  \end{quote}
\item
  \begin{quote}
  исключения лишних копирований при нормализации.
  \end{quote}
\end{itemize}

\textbf{Лабораторная работа №5. Лемматизация / Стемминг и поиск}

Цель работы --- реализовать ранжированный поиск по корпусу (BM25) с
поддержкой стемминга и оценкой качества по набору запросов.

Корпус загружается из MongoDB (docs.text), поэтому индекс строится «в
памяти» без зависимости от файловой структуры. Для идентификации
документов используются пары (source,id) или url; дополнительно
сохраняются title/url для вывода результатов.

Построение данных для поиска выполняется без std::map/unordered\_map:

1) для каждого документа считаются токены (и, при включённой опции,
применён русский стеммер);

2) токены сортируются (merge sort) и сворачиваются в tf по документу;

3) формируется список тройек (term, doc, tf), затем он сортируется по
(term, doc);

4) из отсортированного списка строятся две структуры:

• лексикон (term → start, df) в виде массива LexEntry (термы в
сортированном виде);

• postings --- массив Posting \{doc, tf\} для всех термов подряд.

Поиск: запрос токенизируется (и стеммится при необходимости), для
каждого терма бинарным поиском находится запись в лексиконе и
последовательно обрабатывается соответствующий участок postings с
накоплением BM25-оценки.

Оценка качества: при запуске с флагом -\/-eval читается файл запросов в
формате «query\textless TAB\textgreater comma-separated
relevant-doc-ids», затем считаются P@1, P@3, P@5, P@10 и R@1, R@3, R@5,
R@10.

\textbf{Результаты и выводы (ЛР5)}

Средняя длина документа (в токенах): avgdl=4803.71. Параметры BM25:
k1=1.5, b=0.75, top\_k=10.

\textbf{Оценка объёма индекса (по triples):}

\begin{itemize}
\item
  \begin{quote}
  stemming=OFF:\textbf{\hfill\break
  } для 40 000: ≈ \textbf{45 978 498,2} triples
  \end{quote}
\item
  \begin{quote}
  stemming=ON:\textbf{\hfill\break
  } для 40 000: ≈ \textbf{39 082 835,087} triples
  \end{quote}
\end{itemize}

\textbf{Оценка времени построения (линейная, ``оптимистичная''):}

\begin{itemize}
\item
  \begin{quote}
  stemming=OFF: 0.00366 сек/док → \textbf{≈ 146 сек (2.4 мин)}
  \end{quote}
\item
  \begin{quote}
  stemming=ON: 0.00726 сек/док → \textbf{≈ 290 сек (4.8 мин)}
  \end{quote}
\end{itemize}

\textbf{=== EVAL (1 queries) ===}

\textbf{P@1=1 P@3=0.333333 P@5=0.2 P@10=0.1}

\textbf{R@1=1 R@3=1 R@5=1 R@10=1}

\textbf{\textgreater{} салат с крабовыми палочками}

\textbf{koolinar/149636 9.95319}

\textbf{povarenok/69819 9.77824}

\textbf{koolinar/149402 9.68086}

\textbf{koolinar/149420 9.62964}

\textbf{koolinar/149891 9.55134}

\textbf{koolinar/149389 9.44132}

\textbf{koolinar/149507 9.24544}

\textbf{povarenok/69798 9.12493}

\textbf{povarenok/69970 8.86782}

\textbf{koolinar/149523 8.81013}

Число справа (9.95319, 9.77824, \ldots) - это BM25 score.

Чем больше число --- тем сильнее документ ``похож'' на запрос по модели
BM25.

\textbf{Лабораторная работа №6. Булев индекс}

Цель работы --- построить булев индекс по корпусу, пригодный для
операций AND/OR/NOT, в собственном бинарном формате, а также создать
«прямой» индекс (метаданные документов) для будущей выдачи.

В актуальном пайплайне документы также берутся из MongoDB (docs.text).
Это гарантирует, что индекс строится по тому же корпусу, который собран
роботом в ЛР2.

Формат файла индекса (BIDX) реализован как расширяемый: в начале
записывается заголовок с magic/version и таблица секций. Каждая секция
имеет идентификатор, смещение, размер и auxiliary-поле для будущих
расширений.

Секции индекса:

1) STRS --- строковая таблица (сырой байтовый массив строк).

2) DOCS --- прямой индекс: массив записей DocRec \{docId,
url\_off/url\_len, title\_off/title\_len, text\_bytes, reserved\}. url и
title берутся из MongoDB (url и title).

3) TERMS --- лексикон: массив TermRec \{term\_off/term\_len, post\_off,
df, reserved\}.

4) POSTS --- обратный индекс: для каждого терма хранится список
документов в виде gap-кодированных varint (d1+1, d2-d1, ...), что
уменьшает размер хранения.

Построение: после токенизации каждого документа берутся уникальные термы
в документе (без tf, так как булев индекс), собираются пары (term,
docId), сортируются (merge sort) и сворачиваются по термам, формируя
postings и df.

Результат сохраняется в файл index.bidx (или заданный второй аргумент
командной строки).

\textbf{Результаты и выводы (ЛР6)}

Общий объём текста: 1 309 376.863 КБ. Время построения булева индекса:
150.417 с.

avg\_term\_len\_bytes: 13.9205

На один документ: 3.76 мс/док; на 1 КБ текста: 0.115 мс/КБ.

Сравнение средних длин: средняя длина терма 13.9205 \textgreater{}
средней длины токена 10.8474, т.к. в ЛР6 терм --- это уникальная
нормализованная строка после фильтрации (редко встречаются короткие
служебные слова), а средняя длина токена «разбавляется» короткими
словами и частицами. Кроме того, в корпусе рецептов много составных
наименований ингредиентов и длинных слов.

Внутреннее представление: документы читаются из MongoDB (docs.text,
docs.title, docs.url). После токенизации для каждого документа
формируется список уникальных термов (tf не нужен), затем пары (term,
docId) сортируются и сворачиваются в postings-листы с df. Postings
кодируются gap+varint (экономия места).

Сортировка: используется merge sort (стабильный, предсказуемый O(n log
n), не требует STL контейнеров). Минусы --- дополнительная память под
tmp-массив и цена копирования строк; для больших корпусов выгоднее
хранить термы в string table и сортировать по (hash/off) или применять
external merge.

Масштабирование: при увеличении корпуса в 10× (до 400k документов) при
линейной оценке время вырастет \textasciitilde в 10×, но фактически
основным ограничением станет RAM и пропускная способность памяти/диска
при сортировке. В 100× и 1000× без внешней сортировки/инкрементального
построения индексатор перестанет помещаться в память. Решение ---
блочное построение и слияние, а также переход на более компактные
представления (термы по словарю + postings в файле).

\textbf{Лабораторная работа №7. Булев поиск и веб-сервис}

Цель работы --- реализовать булев поиск по построенному индексу и (при
необходимости) простой веб-сервис выдачи.

ЛР7 работает с бинарным индексом, построенным в ЛР6. Поэтому подключения
к MongoDB для выполнения запросов не требуется: для ответа достаточно
данных из index.bidx (прямой индекс для отображения url/title и обратный
индекс для вычисления результата булевых выражений).

Компоненты: парсер булева выражения (термы и операторы AND/OR/NOT со
скобками), модуль чтения BIDX и исполнитель операций над упорядоченными
списками docId (пересечение/объединение/разность).

При реализации веб-части HTTP-сервер отдаёт HTML-страницу с формой
запроса и списком результатов; данные для выдачи берутся из секции DOCS
(url/title).

\textbf{Результаты и выводы (ЛР7)}

Эксперимент: выполнено 5 запросов

QUERY: московский авиационный институт

hits=0 time\_ms=0.0251

QUERY: (красный \textbar\textbar{} желтый) автомобиль

hits=0 time\_ms=0.024

руки !ноги

QUERY: руки !ноги

hits=2771 time\_ms=0.162201

1. Пирог с яблоками -- кулинарный рецепт \textbar{}
https://www.povarenok.ru/recipes/show/69989

2. Печенье "Очаровашки" -- кулинарный рецепт \textbar{}
https://www.povarenok.ru/recipes/show/69987

3. Куриное филе с грибами и соусом "Бешамель" -- кулинарный рецепт
\textbar{} https://www.povarenok.ru/recipes/show/69957

4. Пирожное "Фантазия" -- кулинарный рецепт \textbar{}
https://www.povarenok.ru/recipes/show/69919

5. Пряничный хлеб -- кулинарный рецепт \textbar{}
https://www.povarenok.ru/recipes/show/69911

6. Потивлеш -- кулинарный рецепт \textbar{}
https://www.povarenok.ru/recipes/show/69906

И т.д\ldots..

QUERY: (курица \&\& !суп) \textbar\textbar{} (котлеты \&\& помидоры)

hits=5752 time\_ms=0.328301

QUERY: !(уха \textbar\textbar{} суп) \&\& (курица \textbar\textbar{}
рыба)

hits=305 time\_ms=0.2205

С ростом коллекции до десятков тысяч документов время запроса в булевом
поиске обычно определяется длинами postings-листов участвующих термов и
сложностью выражения (количество операций AND/OR/NOT). Для редких термов
время почти не растёт, для частых --- растёт пропорционально размеру
списков (в худшем случае близко к O(N)).

Примеры «тяжёлых» запросов: выражения с частыми термами и большим числом
OR, например «(курица \textbar\textbar{} мясо \textbar\textbar{} рыба
\textbar\textbar{} суп \textbar\textbar{} салат) \&\& (лук
\textbar\textbar{} чеснок \textbar\textbar{} перец)». Они заставляют
объединять крупные списки.

Ускорение:

\begin{itemize}
\item
  \begin{quote}
  упорядочивать конъюнкции по возрастанию df (сначала пересекать редкие
  списки),
  \end{quote}
\item
  \begin{quote}
  использовать skip-пойнты в postings,
  \end{quote}
\item
  \begin{quote}
  кэшировать результаты подвыражений.
  \end{quote}
\end{itemize}

Проверка корректности:

\begin{itemize}
\item
  \begin{quote}
  ручные контрольные запросы по очевидным словам и сверка с вхождениями
  в тексте/заголовке,
  \end{quote}
\item
  \begin{quote}
  сравнение булевой выдачи с «наивным» поиском по Python (полный проход
  по docs.text) на небольшой подвыборке,
  \end{quote}
\item
  \begin{quote}
  тесты парсера выражений (скобки, приоритеты, NOT) и крайние случаи
  (пустые результаты, неизвестные термы).
  \end{quote}
\end{itemize}

Можно запустить ./lab7\_server /home/egor/search/6/build/index.bidx 8080

И прямо в поле ввода вводить поисковые запросы

\textbf{Вывод по лабораторным работам}

В ходе выполнения лабораторных работ был последовательно реализован
полный конвейер информационного поиска: от обработки текстового корпуса
до выполнения поисковых запросов различной сложности. На этапе
токенизации и статистического анализа были выработаны правила разбиения
текста, обеспечивающие баланс между качеством выделяемых токенов и
скоростью обработки. Анализ распределения термов подтвердил соответствие
корпуса закону Ципфа, что указывает на его естественную языковую
структуру и пригодность для построения поисковых индексов.

Булев индекс и система логического поиска продемонстрировали высокую
эффективность при обработке точных и составных запросов с логическими
операторами. Выбранное бинарное представление индекса позволило добиться
компактного хранения и быстрого выполнения поисковых операций, а перенос
корпуса и метаданных в базу данных обеспечил масштабируемость и
воспроизводимость экспериментов. В совокупности выполненные работы
показали, что разработанная система способна корректно и эффективно
обрабатывать большие коллекции документов, а архитектурные решения
позволяют масштабировать её на порядки больший объём данных при переходе
к блочной индексации и внешней сортировке.

\end{document}
